\documentclass[14pt,a4paper]{extarticle}
\usepackage{fontspec}
\setmainfont{Times New Roman}
\setmonofont{Consolas}
\usepackage[main=russian,english]{babel}
\usepackage[a4paper,left=30mm,right=15mm,top=20mm,bottom=20mm]{geometry}
\usepackage{setspace}
\usepackage{indentfirst}
\usepackage{longtable,booktabs,array,calc}
\usepackage{amsmath,amssymb}
\usepackage[hidelinks]{hyperref}
\usepackage{xurl}
\setlength{\emergencystretch}{3em}
\onehalfspacing
\setlength{\parindent}{1.25cm}
\setlength{\parskip}{0pt}
\setcounter{secnumdepth}{2}
\providecommand{\tightlist}{\setlength{\itemsep}{0pt}\setlength{\parskip}{0pt}}
\begin{document}
\begin{titlepage}
\centering
\vspace*{3cm}
{\bfseries Курсовая работа\par}
\vspace{1.5cm}
Дисциплина: Разработка профессиональных приложений\par
\vspace{1.5cm}
Тема: Разработка кроссплатформенного приложения для интеллектуального управления задачами и производственными процессами в агропромышленном комплексе\par
\vfill
Москва, 2026\par
\end{titlepage}
\setcounter{page}{2}
\tableofcontents
\clearpage

\section*{Введение}\addcontentsline{toc}{section}{Введение}

Планирование работ в агропромышленном комплексе связано с ограниченными сроками, сезонностью и одновременным выполнением множества задач. В одном проекте могут сочетаться подготовка оборудования, закупка материалов, проведение полевых работ, проверки и оформление результатов. Даже когда каждое действие само по себе несложно, потеря информации о сроке или приоритете одной задачи влияет на последовательность остальных работ. Для организации такого процесса полезно единое приложение, в котором задачи можно связать с проектом, увидеть в календаре, дополнить заметками и распределить по значимости.

Актуальность темы обусловлена необходимостью видеть связанные работы, их сроки и состояние в одном месте. При раздельном ведении списка задач, расписания и рабочих записей пользователю приходится сопоставлять сведения вручную. Кроссплатформенное приложение позволяет обращаться к единому набору данных с компьютера и телефона, а явные правила приоритизации помогают выделять дела, требующие внимания в первую очередь.

В данной курсовой работе рассматривается разработанный программный прототип WorkNotes. Его назначение шире одного хозяйства: приложение представляет собой универсальный планировщик проектов, задач, заметок и расписания. Агропромышленный комплекс выбран как предметная область для демонстрации сценария применения. Проект «Подготовка к посевной», например, может объединять проверку техники, получение материалов и контроль дат выполнения. При этом программа не содержит специализированных модулей агрономической диагностики, учёта урожайности или управления сельскохозяйственной техникой. Такое разграничение необходимо для точного описания фактически реализованного результата.

Под интеллектуальным управлением в рамках работы понимается \textbf{формализованная поддержка приоритизации}. Система рассчитывает положение задачи в матрице Эйзенхауэра по её приоритету и сроку, допускает ручную корректировку срочности и показывает результат в четырёх секторах. Дополнительно метод MoSCoW позволяет обозначить обязательные, важные, желательные и отложенные задачи проекта. Эти два механизма отвечают на разные вопросы: матрица помогает определить очередность внимания к текущим делам, а MoSCoW --- согласовать состав работ на выбранном этапе. В приложении не используются машинное обучение и прогнозирование, поэтому термин «интеллектуальное» относится к явным правилам классификации и удобному представлению решений.

\textbf{Цель работы} --- разработать и описать кроссплатформенное приложение, которое помогает структурировать проектные задачи, сроки и связанные материалы, а также поддерживает приоритизацию на основе матрицы Эйзенхауэра и MoSCoW.

Для достижения цели поставлены следующие задачи:
\begin{enumerate}
\item проанализировать пользовательские сценарии и сформулировать требования к приложению;
\item выбрать технологии для клиентской и серверной частей;
\item спроектировать модель данных и HTTP API;
\item реализовать интерфейсы проектов, задач, календаря, расписания и заметок;
\item реализовать правила матрицы Эйзенхауэра и группировку MoSCoW;
\item проверить классификацию задач, работу основных операций и сохранение данных.
\end{enumerate}

Объектом рассмотрения являются процессы планирования и контроля задач. Предметом --- программные методы их хранения, отображения и приоритизации в кроссплатформенном приложении.

В работе использованы анализ пользовательских сценариев и требований, проектирование структуры данных и HTTP API, разработка программного прототипа и функциональное тестирование серверной логики.

Практическая значимость работы заключается в том, что прототип объединяет задачи, календарь, расписание и связанные заметки в одном приложении. Пользователь может просматривать работы в нескольких представлениях и выбирать порядок действий с помощью матрицы Эйзенхауэра и MoSCoW. Практический результат представляет собой исходный код Flutter-клиента, сервера FastAPI и базы SQLite, а также набор автоматизированных проверок серверной логики. Клиентский код ориентирован на Windows и Android; подтверждённые в этой работе проверки относятся прежде всего к серверной логике и статическому анализу Dart.

\section{Анализ предметной области и постановка задачи}

\subsection{Сценарий планирования работ}

В ходе подготовки к сезонным работам ответственному сотруднику требуется видеть не только перечень поручений. Для каждой задачи важны содержание, дата начала, конечный срок, текущее состояние и принадлежность к проекту. Описание может включать конкретные действия, а сопутствующая заметка --- результаты осмотра или условия выполнения. Отдельные мероприятия проводятся по расписанию и повторяются еженедельно либо раз в две недели. Если сведения распределены между перепиской, календарём и бумажными записями, пользователю приходится вручную сопоставлять их при каждом изменении плана.

Для демонстрации выбран проект «Подготовка к посевной». В нём можно создать задачи «Проверить состояние сеялки», «Подготовить план закупки семян», «Уточнить готовность участка», «Оформить итоговый отчёт». К первой задаче прикрепляется заметка с результатом проверки; дата контрольного осмотра указывается в расписании. У каждой задачи есть собственные срок и приоритет. Пользователь получает представление о работах в списке, на доске по состоянию, в календаре, в матрице Эйзенхауэра и по категориям MoSCoW. Это один учебный сценарий применения универсального приложения, а не утверждение о внедрении в конкретное хозяйство.

Планирование целесообразно разделить на три уровня. Проект группирует связанные задачи. Задача описывает единицу работы и её состояние. Расписание фиксирует событие во времени: например, контрольное совещание или осмотр. Заметка хранит пояснение и может быть связана с одной или несколькими задачами. Такое разделение предотвращает смешение неизменного контекста, действия и повторяющегося события. В частности, перенос даты встречи не должен автоматически изменять срок задачи, а изменение статуса задачи не обязано удалять относящиеся к ней записи.

Доступность приложения на разных устройствах повышает удобство: работа с подробным списком и заметками возможна за компьютером, а быстрый просмотр сроков и изменение статуса --- на телефоне. В рассматриваемой архитектуре данные хранятся на сервере. Поэтому для работы клиента требуется доступ к API; наличие Flutter-кода для Android само по себе не означает подтверждённую работу всей системы в автономном режиме или с удалённым сервером в производственной сети.

\subsection{Требования к функциональности}

Основной пользовательский объект --- задача. При создании указываются заголовок, необязательное описание, проект, статус, дата начала, необязательный крайний срок и приоритет. Дополнительно можно задать теги, ручную срочность и категорию MoSCoW. Значения статуса ограничены состояниями «Сделать», «В работе» и «Готово». Приоритет принимает значения low, medium и high. Срок выполнения не может быть раньше даты начала. Эти ограничения вводятся в модели запроса и дополнительно контролируются при изменении существующей записи.

Система должна поддерживать создание, просмотр, изменение и удаление задачи, поиск и фильтрацию, сортировку по сроку или приоритету, выделение просроченных задач и календарное отображение. Проект создаётся отдельно; задача может не принадлежать ни одному проекту. Фильтр проекта помогает рассматривать работы одного направления без смешения с остальными. Теги позволяют выделять группы вроде «техника», «закупки» и «документы», но сами по себе не устанавливают порядок выполнения.

Матрица Эйзенхауэра предназначена для активных задач. Система делит их на четыре группы: срочные и важные; несрочные и важные; срочные и менее важные; несрочные и менее важные. Важность выводится из высокого приоритета. Срочность определяется сроком --- сегодняшний день либо ближайшие два дня, включая просроченный срок, --- если пользователь не задал её вручную. Завершённые задачи не включаются в матрицу, но остаются в обычном списке. При перемещении задачи между секторами изменяются поля приоритета и срочности, которые определяют новое положение.

MoSCoW решает иную задачу. Категория Must обозначает обязательную работу текущего этапа, Should --- важную работу, Could --- желательную, Won't --- отложенную. Предусмотрена также группа без категории. Присвоение категории не меняет статус задачи, её обычный приоритет или положение в матрице. В интерфейсе группировка наиболее полезна при выборе проекта: это позволяет оценить его состав, не смешивая задачи разных направлений. Метод применяется к перечню работ; в прототипе категории назначаются вручную и доступны без механизма согласования между несколькими пользователями.

Заметки должны хранить заголовок и форматированный текст. Для текста предусмотрены абзацы, заголовки, маркированные и нумерованные элементы, а также оформление фрагментов. Заметку можно связать с задачами, на её основе создать задачу и приложить изображение либо аудиофайл допустимого формата. Расписание должно поддерживать одноразовую запись и повторение раз в одну или две недели, ограничение даты окончания, отмену одного события в серии. Календарь по запросу объединяет задачи со сроком и экземпляры событий расписания.

\begingroup\small
\begin{longtable}[]{@{}
  >{\raggedright\arraybackslash}p{(\linewidth - 4\tabcolsep) * \real{0.3333}}
  >{\raggedright\arraybackslash}p{(\linewidth - 4\tabcolsep) * \real{0.3333}}
  >{\raggedright\arraybackslash}p{(\linewidth - 4\tabcolsep) * \real{0.3333}}@{}}
\toprule\noalign{}
\begin{minipage}[b]{\linewidth}\raggedright
\textbf{Группа требований}
\end{minipage} & \begin{minipage}[b]{\linewidth}\raggedright
\textbf{Реализуемое действие}
\end{minipage} & \begin{minipage}[b]{\linewidth}\raggedright
\textbf{Проверяемый результат}
\end{minipage} \\
\midrule\noalign{}
\endhead
\bottomrule\noalign{}
\endlastfoot
Проекты & Создать проект и выбрать его & Список задач ограничен проектом \\
Задачи & Создать, изменить, завершить, удалить & Изменения сохраняются в SQLite \\
Сроки & Установить дату и открыть календарь & Задача видна на соответствующий день \\
Матрица & Определить важность и срочность & Активная задача попадает в один сектор \\
MoSCoW & Назначить категорию & Активная задача видна в нужной группе \\
Заметки & Связать текст с задачей & Связь доступна при повторном чтении \\
Расписание & Задать повторение и исключение & Отменённое событие пропадает из выбранной даты \\
\end{longtable}
\endgroup

Таблица 1 --- Основные функции и ожидаемые результаты

\subsection{Нефункциональные требования и ограничения}

Приложение должно сохранять данные после перезапуска сервера, возвращать понятные ошибки на некорректный запрос и обеспечивать одинаковую семантику задач в разных видах интерфейса. Для календаря особенно важно, чтобы изменение даты не приводило к расхождению между списком задач и дневным представлением. Логика матрицы должна быть воспроизводимой: по тем же приоритету, сроку и ручной отметке пользователь получает тот же сектор в указанный день. При наступлении новых суток вычисляемая срочность может измениться без редактирования записи; интерфейс обновляет данные около полуночи.

Серверная часть разделена по предметным модулям, чтобы изменение правил расписания не требовало переписывать обработчики заметок. В базе используются первичные и внешние ключи для сохранения связности записей. При удалении связанной сущности применяются заданные правила для зависимых строк. Для примера, теги удалённой задачи не должны оставаться самостоятельными записями, а удаление проекта не уничтожает автоматически саму задачу: её связь с проектом снимается.

Текущий прототип не содержит учётных записей, участников команды, назначения исполнителя, разграничения доступа, зависимостей задач и учёта ресурсов хозяйства. Соответственно, «управление производственными процессами» в названии работы понимается как планирование набора работ и контроль сроков в проекте, а не как полный производственный контур предприятия. При эксплуатации вне локальной сети потребуются отдельные решения по аутентификации, правам, резервному копированию и размещению сервера. Эти вопросы обозначены как развитие системы и не выдаются за реализованные возможности.

\subsection{Выбор технологий}

Клиентская часть написана на Dart с использованием Flutter. Один набор экранов применяется для настольного и мобильного интерфейса; компоновка матрицы адаптируется к ширине: на широком экране четыре сектора образуют квадрат, на узком располагаются вертикально. Flutter поддерживает сборку приложений для мобильных и настольных платформ, что соответствует выбранной цели {[}1{]}. Наличие общей кодовой базы уменьшает необходимость отдельно реализовывать логику экранов на каждом языке платформы, хотя сборку и поведение каждой целевой платформы нужно проверять отдельно.

Сервер разработан на Python с FastAPI. Маршруты по задачам, проектам, заметкам и расписанию подключаются к одному приложению через APIRouter. Такой подход позволяет разнести крупный API по файлам и сохранять общую точку запуска {[}2{]}. Pydantic-модели описывают входящие значения и структуру ответов; FastAPI предоставляет автоматически формируемую документацию API {[}3{]}. На клиенте HTTP-запросы собраны в отдельном ApiClient, а компоненты интерфейса разделены по функциональным каталогам.

Для хранения выбран SQLite. База расположена в одном файле на сервере, что упрощает локальную разработку и демонстрацию. Схема содержит связи между проектами, задачами, тегами, заметками, приложениями к ним и событиями расписания. Контроль внешних ключей включается на соединениях, для которых он необходим: в SQLite это явное требование {[}4{]}. Выбор SQLite соответствует масштабу учебного прототипа; при расширении до работы нескольких хозяйств и большого числа пользователей необходимо отдельно оценить нагрузку, права и схему развертывания.

\section{Проектирование структуры и методов приоритизации}

\subsection{Общая архитектура}

Система состоит из клиента Flutter, сервера FastAPI и файла базы SQLite. Клиент показывает экраны и принимает действия пользователя, API проверяет запросы и вызывает предметную логику, репозитории выполняют SQL-операции. Результаты возвращаются в JSON. Изображения и аудиозаписи заметок сохраняются отдельными файлами в каталоге загрузок; в таблице хранится их метаинформация и связь с заметкой. Файл базы и каталог вложений должны сохраняться вместе при резервном копировании.

\begingroup\small
\begin{longtable}[]{@{}
  >{\raggedright\arraybackslash}p{(\linewidth - 4\tabcolsep) * \real{0.3333}}
  >{\raggedright\arraybackslash}p{(\linewidth - 4\tabcolsep) * \real{0.3333}}
  >{\raggedright\arraybackslash}p{(\linewidth - 4\tabcolsep) * \real{0.3333}}@{}}
\toprule\noalign{}
\begin{minipage}[b]{\linewidth}\raggedright
\textbf{Уровень}
\end{minipage} & \begin{minipage}[b]{\linewidth}\raggedright
\textbf{Основные компоненты}
\end{minipage} & \begin{minipage}[b]{\linewidth}\raggedright
\textbf{Ответственность}
\end{minipage} \\
\midrule\noalign{}
\endhead
\bottomrule\noalign{}
\endlastfoot
Интерфейс & \path{app/lib/features}, \path{app/lib/core/api_client.dart} & Экран, форма, вызов API, состояние загрузки \\
HTTP API & \path{backend/main.py}, модули \texttt{router.py} & Маршруты, коды ответа, проверка входа \\
Сервис & Модули \texttt{service.py} & Правила задач, календаря, связей и повторов \\
Доступ к данным & Модули \texttt{repository.py} & SQL-запросы и преобразование строк \\
Хранение & \path{backend/tasks.db}, \path{backend/uploads} & Долговременные данные и вложения \\
\end{longtable}
\endgroup

Таблица 2 --- Уровни приложения

Например, при создании задачи форма формирует поля и отправляет POST-запрос. Схема TaskCreate проверяет типы, допустимые значения и согласованность дат. Роутер передаёт модель в сервис; сервис проверяет существование выбранного проекта; репозиторий вставляет строку и теги. В ответ клиент получает задачу с идентификатором и вычисляемыми полями. Этот маршрут показывает разделение ответственности: интерфейс не должен самостоятельно записывать данные в SQLite, а SQL-слой не определяет внешний вид страницы.

В \path{backend/main.py} создаётся объект FastAPI, инициализируется база и подключаются четыре маршрутизатора. Наличие \path{/health} позволяет проверить, отвечает ли процесс сервера. Команда локального запуска использует Uvicorn и адрес \texttt{127.0.0.1:8000}. При таком адресе приложение на том же компьютере может обращаться к API; Android-устройству потребуется адрес сервера, доступный ему по сети. Сам факт сборки Android-клиента не гарантирует соединение с локальным \texttt{127.0.0.1} компьютера.

Клиент содержит модели Task, Project, CalendarMonth, ScheduleEntry и Note. Экраны задач и вспомогательные диалоги располагаются в \path{features/tasks}, заметки --- в \path{features/notes}, расписание --- в \path{features/schedule}. \texttt{ApiClient} инкапсулирует HTTP-операции, поэтому компоненту матрицы не требуется знать, как формируется URL. После изменения задачи экран загружает актуальные данные; это особенно важно, когда сектор вычисляется по сроку на сервере.

\subsection{Модель данных}

Центральная сущность \texttt{tasks} содержит идентификатор, заголовок, описание, состояние, ссылку на проект, даты, обычный приоритет, ручную срочность и категорию MoSCoW. Поле \texttt{urgent} допускает три состояния: \texttt{1} --- срочно, \texttt{0} --- не срочно, \texttt{NULL} --- вычислить по дате. Поле \texttt{moscow} допускает одно из четырёх кодовых значений или \texttt{NULL}, если решение ещё не принято. Трёхзначная модель срочности позволяет явно сохранить выбор пользователя и одновременно оставить автоматическое правило по умолчанию. Категория MoSCoW не дублирует приоритет: важная для текущего этапа задача может иметь любой срок.

\begingroup\small
\begin{longtable}[]{@{}
  >{\raggedright\arraybackslash}p{(\linewidth - 4\tabcolsep) * \real{0.3333}}
  >{\raggedright\arraybackslash}p{(\linewidth - 4\tabcolsep) * \real{0.3333}}
  >{\raggedright\arraybackslash}p{(\linewidth - 4\tabcolsep) * \real{0.3333}}@{}}
\toprule\noalign{}
\begin{minipage}[b]{\linewidth}\raggedright
\textbf{Таблица}
\end{minipage} & \begin{minipage}[b]{\linewidth}\raggedright
\textbf{Существенные поля}
\end{minipage} & \begin{minipage}[b]{\linewidth}\raggedright
\textbf{Назначение связи}
\end{minipage} \\
\midrule\noalign{}
\endhead
\bottomrule\noalign{}
\endlastfoot
\texttt{projects} & \texttt{id}, \texttt{name} & Группировка задач и записей расписания \\
\texttt{tasks} & \texttt{id}, \texttt{project\_id}, \texttt{status}, \texttt{deadline}, \texttt{priority}, \texttt{urgent}, \texttt{moscow} & Основная единица работы \\
\texttt{task\_tags} & \texttt{task\_id}, \texttt{name}, \texttt{name\_key}, \texttt{position} & Несколько тегов одной задачи \\
\texttt{notes} & \texttt{id}, \texttt{title}, \texttt{content}, даты изменения & Текстовая информация \\
\texttt{note\_tasks} & \texttt{note\_id}, \texttt{task\_id} & Связь многие ко многим \\
\texttt{note\_attachments} & \texttt{id}, \texttt{note\_id}, тип, размер, имя файла & Метаданные вложения \\
\texttt{schedule\_entries} & \texttt{id}, \texttt{project\_id}, \texttt{first\_date}, время, период & Шаблон события \\
\texttt{schedule\_exceptions} & \texttt{entry\_id}, \texttt{date} & Отмена отдельного повторения \\
\end{longtable}
\endgroup

Таблица 3 --- Основные таблицы базы данных

Связь проекта с задачей необязательна. При удалении проекта идентификатор проекта в задаче становится пустым, и сама задача остаётся в базе. Удаление задачи каскадно удаляет её теги и строки связи с заметками. Это различие важно: проект служит группой, но не является единственным владельцем содержимого задачи. У расписания есть собственный необязательный \texttt{project\_id}, поэтому события могут быть как общими, так и относящимися к одному проекту.

Заметки хранят форматированные блоки как JSON в текстовом поле. Отдельная таблица \texttt{note\_tasks} нужна, поскольку одна заметка может сопровождать несколько задач, а одна задача --- иметь несколько заметок. Файловые вложения не записываются в основной текст заметки: поле \texttt{storage\_name} связывает метаданные с файлом, имя которого назначает сервер. При удалении заметки необходимо учитывать и строки метаданных, и физические файлы.

Расписание хранит не каждый экземпляр повторяющегося события, а запись-шаблон с начальной датой, временем и периодом \texttt{repeat\_every\_weeks}: ноль, одна или две недели. При запросе диапазона сервис рассчитывает конкретные даты. Исключения помещаются в отдельную таблицу и позволяют отменить одну дату, не уничтожая всю серию. Например, чередование занятий или контрольных мероприятий «неделя А --- неделя Б» реализуется двумя записями с шагом две недели и начальными датами, отличающимися на семь дней.

Создание недостающих таблиц и добавление полей к существующей базе выполняются в \texttt{init\_db()}. Для задач, созданных до появления \texttt{urgent} и \texttt{moscow}, новые значения остаются \texttt{NULL}. Это сохраняет старые записи и даёт им предсказуемое поведение: срочность рассчитывается по сроку, категория MoSCoW остаётся неназначенной. Повторный вызов инициализации не должен очищать данные. Такая миграция подходит для текущей стадии прототипа; при большом количестве изменений схемы потребовалась бы отдельная система версионирования миграций.

\subsection{Матрица Эйзенхауэра как правило классификации}

Матрица использует две независимые оси: важность и срочность. В исходной задаче нет отдельного поля важности; по принятому в проекте правилу важной считается задача с приоритетом \texttt{high}. Приоритеты \texttt{medium} и \texttt{low} для цели двоичной классификации образуют группу менее важных. Пользователь может изменить приоритет в редакторе либо переместить задачу в сектор другой важности. Такой выбор упрощает форму, но требует пояснения: «менее важная» в матрице не означает, что задача вообще не нужна.

Для срочности сначала проверяется поле \texttt{urgent}. Если там явно указано \texttt{true} или \texttt{false}, именно оно определяет сектор. При \texttt{NULL} задача срочная, когда у неё есть крайний срок не позже чем через два календарных дня от текущей даты. Отсутствие срока означает отсутствие вычисляемой срочности. Просроченные активные задачи удовлетворяют условию близкого срока, поскольку их дата уже прошла. Отметка просрочки вычисляется отдельно: она истинна, когда срок раньше сегодняшнего дня и статус не равен \texttt{done}.

Обозначим важность через I, а срочность через U. Тогда I = 1 при \texttt{priority\ =\ high}, иначе I = 0. Если \texttt{urgent} задано, U равно его значению; иначе U = 1 при наличии срока и условии \texttt{deadline} $\leq$ \texttt{today + 2 дня}. Четыре возможные пары \texttt{(U,\ I)} однозначно определяют сектор. На границе суток вычисление может изменить U без изменения строки в базе, поэтому его не следует хранить как постоянно обновляемый отдельный флаг.

\begingroup\small
\begin{longtable}[]{@{}
  >{\raggedright\arraybackslash}p{(\linewidth - 6\tabcolsep) * \real{0.2857}}
  >{\raggedright\arraybackslash}p{(\linewidth - 6\tabcolsep) * \real{0.2857}}
  >{\raggedright\arraybackslash}p{(\linewidth - 6\tabcolsep) * \real{0.2143}}
  >{\raggedright\arraybackslash}p{(\linewidth - 6\tabcolsep) * \real{0.2143}}@{}}
\toprule\noalign{}
\begin{minipage}[b]{\linewidth}\raggedright
\textbf{Срочность U}
\end{minipage} & \begin{minipage}[b]{\linewidth}\raggedright
\textbf{Важность I}
\end{minipage} & \begin{minipage}[b]{\linewidth}\raggedright
\textbf{Сектор}
\end{minipage} & \begin{minipage}[b]{\linewidth}\raggedright
\textbf{Пример в учебном проекте}
\end{minipage} \\
\midrule\noalign{}
\endhead
\bottomrule\noalign{}
\endlastfoot
1 & 1 & Срочно и важно & Устранить выявленную неисправность перед работой \\
0 & 1 & Важно, но не срочно & Подготовить план закупок на следующий этап \\
1 & 0 & Срочно, но менее важно & Уточнить время текущего созвона \\
0 & 0 & Не срочно и менее важно & Упорядочить старые черновые заметки \\
\end{longtable}
\endgroup

Таблица 4 --- Правила и учебные примеры матрицы

Классификация выполняется на сервере для отдельной задачи и для ответа \path{GET\ /tasks/matrix}. В групповой ответ включаются только задачи, у которых статус отличен от \texttt{done}; можно передать \texttt{project\_id}. Если клиент перемещает задачу в другой сектор, он должен обновить соответствующие исходные поля. Например, перевод в сектор «срочно и важно» задаёт высокую важность и ручную срочность. Перевод между секторами одинаковой важности не должен произвольно менять низкий и средний уровни приоритета. После записи сервер заново вычисляет сектор, а клиент обновляет экран.

Алгоритм прозрачен для проверки и защиты: любому результату соответствует набор входных полей и явное правило. Он помогает пользователю быстрее увидеть сочетание срока и значимости, но не измеряет экономический эффект задачи, не прогнозирует продолжительность работ и не заменяет решение руководителя. Для задач без дедлайна или при исключительных обстоятельствах пользователь может задать срочность вручную.

\subsection{Группировка MoSCoW и её отличие от матрицы}

MoSCoW применяется для определения состава работы на выбранном этапе проекта. Метод различает Must, Should, Could и Won't this time; практический смысл категорий связан с необходимостью результата и ограничением времени {[}5{]}. В данной реализации значение присваивается каждой задаче вручную. \texttt{Won’t} означает отложить задачу на рассматриваемом этапе, а не удалить её и не перевести автоматически в статус «Готово». Пятая техническая группа «Без категории» помогает обнаружить задачи, по которым решение не принято.

На примере проекта «Подготовка к посевной» обязательной работой можно обозначить проверку исправности оборудования, важной --- подготовку резервного перечня материалов, желательной --- оформление дополнительной памятки, отложенной --- обновление несрочных иллюстраций отчёта. Это демонстрационные назначения, сделанные автором сценария, а не экспертная агрономическая оценка конкретного хозяйства. В отличие от матрицы, MoSCoW не вычисляется по дедлайну и не изменяется автоматически с наступлением новой даты.

\begingroup\small
\begin{longtable}[]{@{}
  >{\raggedright\arraybackslash}p{(\linewidth - 4\tabcolsep) * \real{0.3333}}
  >{\raggedright\arraybackslash}p{(\linewidth - 4\tabcolsep) * \real{0.3333}}
  >{\raggedright\arraybackslash}p{(\linewidth - 4\tabcolsep) * \real{0.3333}}@{}}
\toprule\noalign{}
\begin{minipage}[b]{\linewidth}\raggedright
\textbf{Признак}
\end{minipage} & \begin{minipage}[b]{\linewidth}\raggedright
\textbf{Матрица Эйзенхауэра}
\end{minipage} & \begin{minipage}[b]{\linewidth}\raggedright
\textbf{MoSCoW}
\end{minipage} \\
\midrule\noalign{}
\endhead
\bottomrule\noalign{}
\endlastfoot
Вопрос & На что обратить внимание сейчас? & Что включить в текущий объём работ? \\
Входные данные & Приоритет, срок, ручная срочность & Отдельная категория, заданная человеком \\
Группы & Четыре сочетания важности и срочности & Must, Should, Could, Won't, без категории \\
Изменение со временем & Возможно при приближении дедлайна & Только при решении пользователя \\
Влияние на статус & Отсутствует & Отсутствует \\
\end{longtable}
\endgroup

Таблица 5 --- Различие двух механизмов приоритизации

Одна и та же задача может одновременно быть \texttt{Must} в MoSCoW и находиться в секторе «важно, но не срочно» матрицы. Противоречия нет: первое решение касается необходимости работы в проекте, второе --- времени внимания к ней. Независимое хранение \texttt{moscow} и \texttt{urgent} предотвращает непреднамеренное изменение одной методики при работе с другой. При удалении проекта задача не теряет категорию, хотя перестаёт относиться к этому проекту; это следует учитывать при последующей ручной классификации.

\subsection{Проектирование API и ошибок}

HTTP API использует ресурсы \path{/projects}, \path{/tasks}, \path{/notes} и \path{/schedule}. Операции чтения выполняются GET, создание --- POST, частичное изменение --- PATCH, удаление --- DELETE. Для связей заметки и задачи предусмотрены PUT и DELETE. Выделенные маршруты \path{/tasks/matrix} и \path{/tasks/moscow} возвращают сгруппированные активные задачи; маршрут \path{/tasks/calendar} --- записи по дням заданного месяца. Структура ответа определяется моделями Pydantic, поэтому интерфейс может работать с предсказуемыми полями.

\begingroup\small
\begin{longtable}[]{@{}
  >{\raggedright\arraybackslash}p{(\linewidth - 4\tabcolsep) * \real{0.3333}}
  >{\raggedright\arraybackslash}p{(\linewidth - 4\tabcolsep) * \real{0.3333}}
  >{\raggedright\arraybackslash}p{(\linewidth - 4\tabcolsep) * \real{0.3333}}@{}}
\toprule\noalign{}
\begin{minipage}[b]{\linewidth}\raggedright
\textbf{Запрос}
\end{minipage} & \begin{minipage}[b]{\linewidth}\raggedright
\textbf{Назначение}
\end{minipage} & \begin{minipage}[b]{\linewidth}\raggedright
\textbf{Существенный параметр или результат}
\end{minipage} \\
\midrule\noalign{}
\endhead
\bottomrule\noalign{}
\endlastfoot
\path{GET\ /tasks} & Список и фильтры & \texttt{project\_id}, статус, тег, приоритет, сортировка \\
\path{GET\ /tasks/matrix} & Четыре сектора & Необязательный \texttt{project\_id} \\
\path{GET\ /tasks/moscow} & Категории MoSCoW & Пять массивов, включая \texttt{unassigned} \\
\path{GET\ /tasks/calendar} & Календарь месяца & Год, месяц, проект, \texttt{include\_schedule} \\
\path{POST\ /tasks} & Создание & Заголовок, даты, статус, приоритет, категории \\
\texttt{PATCH\ /tasks/\{id\}} & Частичное изменение & Только переданные поля \\
\path{GET\ /schedule/occurrences} & Конкретные события & Интервал дат и проект \\
\texttt{POST\ /notes/\{id\}/create-task} & Задача из заметки & Даты, статус и приоритет новой задачи \\
\end{longtable}
\endgroup

Таблица 6 --- Выбранные операции API

Некорректный порядок дат даёт ошибку валидации, отсутствие задачи или проекта --- ответ 404. Создание обычно возвращает 201, успешное удаление --- 204. Ограничения на значения статусов, приоритетов и категории задаются типами и схемой базы. Для обновления необходимо различать «поле не передано» и «передано \texttt{null}»: в первом случае значение сохраняется, во втором для допустимых полей сбрасывается. Например, \texttt{PATCH} с \texttt{urgent:\ null} возвращает автоматическое правило по сроку, а \texttt{moscow:\ null} убирает ручную категорию.

Отдельное значение имеет фильтр \texttt{project\_id}. Для задач он позволяет сопоставлять виды интерфейса в одном и том же контексте. Для расписания сервис строит экземпляры событий на заданном интервале, не храня копию каждого повторения. Для календаря флаг \texttt{include\_schedule=true} добавляет эти события к задачам; без него ответ содержит только задачи. Сохранение необязательного параметра не меняет смысл старого запроса и упрощает развитие клиента.

\section{Реализация приложения и проверка результата}

\subsection{Серверные модули задач и проектов}

В серверной части модули расположены по сущностям. \path{backend/projects} содержит схемы, маршруты, сервис и репозиторий проектов; аналогичная организация используется для \path{backend/tasks}, \path{backend/notes} и \path{backend/schedule}. Такой порядок позволяет проследить путь запроса и локализовать изменения. Например, добавление категории MoSCoW затронуло поле модели задачи, колонку базы, чтение и запись в репозитории, маршрут группировки и клиентский экран. В то же время маршруты заметок не потребовали изменения своей основной структуры.

Проект хранит идентификатор и название. При создании задачи сервис проверяет существование проекта, если его идентификатор передан. Несуществующая ссылка приводит к ответу 404, а отсутствие ссылки допускается. Удаление проекта выполняется отдельно от удаления задач. Такое поведение удобно при закрытии организационной группы: пользователь не теряет накопленные поручения, хотя должен решить, куда перенести задачи без проекта.

В \texttt{TaskCreate} указаны допустимые значения статуса, приоритета и категории, а также ограничение длины заголовка. Валидатор дат запрещает крайний срок раньше даты начала. При PATCH используется модель \texttt{TaskUpdate}, в которой поля необязательны; сервис объединяет переданные изменения с текущей записью и снова проверяет порядок дат. Это защищает от случая, когда новая дата начала становится позже уже сохранённого дедлайна. Заголовок, теги и другие значения также проверяются до передачи в слой SQL.

Теги хранятся отдельно от основной строки задачи. При обработке удаляются лишние пробелы и отклоняются пустые, слишком длинные или повторяющиеся без учёта регистра названия. Благодаря нормализованному ключу \texttt{name\_key} фильтрация по тегу не зависит от того, как пользователь написал буквы. Порядковый номер тега сохраняет его положение в исходном списке. Это пример требования, не видимого на главном экране, но влияющего на качество данных.

Для просрочки используется вычисляемое поле \texttt{is\_overdue}: статус не \texttt{done}, дедлайн существует и он меньше текущей даты. Просроченная задача не переводится автоматически в новое состояние; пользователь может завершить её, изменить срок или удалить. Сортировка по сроку и приоритету дополняет фильтры, но не меняет сохранённые данные. На уровне представления это даёт несколько взглядов на один набор задач без создания их копий.

Запуск сервера выполняется командой \texttt{python -m uvicorn backend.main:app --host 127.0.0.1 --port 8000 --reload} из корня проекта при активном окружении. Страница \path{/docs} содержит интерактивное описание маршрутов, а \path{/health} возвращает ответ \texttt{\{"status":"ok"\}}. Параметр \path{-\/-reload} применяется для локальной разработки; при постоянном размещении параметры запуска и доступ по сети необходимо настроить отдельно. База \path{backend/tasks.db} создаётся при инициализации, причём тесты используют временные файлы, чтобы не менять пользовательские записи.

\subsection{Реализация матрицы Эйзенхауэра}

В серверной модели задачи вычисляемые свойства \texttt{is\_urgent} и \texttt{matrix\_quadrant} возвращаются вместе с обычными полями. Первое свойство проверяет ручное значение \texttt{urgent}, а при его отсутствии сравнивает дедлайн с текущей датой плюс два дня. Второе комбинирует срочность и важность, определённую как \texttt{priority\ ==\ high}. Метод \texttt{task\_matrix} получает список задач, исключает \texttt{done} и помещает каждую активную задачу ровно в одну из четырёх групп. Ответ содержит массивы \texttt{urgent\_important}, \texttt{not\_urgent\_important}, \texttt{urgent\_not\_important} и \texttt{not\_urgent\_not\_important}.

На стороне Flutter компонент \texttt{TaskMatrixBoard} работает с уже загруженными задачами и их вычисленным сектором. В верхней части показано правило классификации; внутри каждого сектора --- название, короткое пояснение, число задач и карточки. При ширине меньше 700 логических пикселей четыре сектора идут один под другим. При большей ширине используются два ряда по два сектора. Таким образом, читаемость сохраняется и на узком мобильном экране, и в широком окне настольного приложения.

На рисунке~\ref{fig:matrix-all} показана матрица в работающем приложении для Windows. Отдельный снимок на рисунке~\ref{fig:matrix-project} демонстрирует те же четыре сектора после выбора проекта «курсовая»: задачи других проектов в нём не отображаются. Снимки сделаны на компьютере; экран Android будет документироваться после проверки мобильной сборки.

\begin{figure}[!htbp]
\centering
\XeTeXpicfile figures/windows_matrix_all.png width \linewidth
\caption{Матрица Эйзенхауэра для всех задач в приложении Windows}
\label{fig:matrix-all}
\end{figure}

\begin{figure}[!htbp]
\centering
\XeTeXpicfile figures/windows_matrix_project.png width \linewidth
\caption{Матрица Эйзенхауэра в проекте «курсовая» в приложении Windows}
\label{fig:matrix-project}
\end{figure}

Команда «Переместить» позволяет выбрать другой сектор. Обработчик в \texttt{TasksPage} преобразует выбор сектора в пару признаков \texttt{important} и \texttt{urgent}, а \texttt{ApiClient} отправляет изменение задачи. Если новая важность совпадает с прежней, текущий средний или низкий приоритет можно сохранить; если она меняется, записывается значение, соответствующее выбранной группе. Срочность после ручного перемещения получает явное значение, поэтому позднейшее изменение даты само по себе не отменит выбор пользователя. Вернуться к автоматическому правилу можно через редактирование задачи и сброс поля.

При работе с матрицей нужно различать три времени: дату начала, крайний срок и текущую дату. Дата начала применяется при проверке допустимого интервала, но сама по себе не определяет срочность. Дедлайн сравнивается с датой на сервере при формировании ответа. Клиент планирует обновление после полуночи, чтобы уже открытый экран не показывал устаревшее вычисленное значение бесконечно долго. В распределённом размещении сервер и клиент должны использовать согласованный календарный день; отдельная настройка часового пояса станет важной при использовании за пределами локального прототипа.

Проверочный пример показывает значение ручного исключения. Задача с высоким приоритетом и сроком завтра автоматически окажется в секторе «срочно и важно». Если пользователь зафиксировал \texttt{urgent\ =\ false}, она перейдёт в «важно, но не срочно» независимо от близкого срока. Если затем передать \texttt{urgent\ =\ null}, автоматическое правило снова сделает задачу срочной. Такая последовательность проверяет не только четыре сектора, но и приоритет ручного решения над вычисленным.

\subsection{Реализация MoSCoW}

Категория MoSCoW хранится в отдельном поле \texttt{moscow} задачи. Допустимы \texttt{must}, \texttt{should}, \texttt{could}, \texttt{wont} и пустое значение. Это поле добавляется к старой базе без принудительного заполнения: ранее созданные задачи попадают в группу \texttt{unassigned}. Метод \texttt{task\_moscow} выбирает активные задачи указанного проекта и распределяет их по пяти массивам. Внутри группы сохраняется сортировка по обычному приоритету, затем по сроку согласно правилам списка. Категория при этом не подменяет исходный приоритет.

В форме задачи пользователь может выбрать категорию одновременно с остальными полями. Альтернативный путь --- открыть вкладку MoSCoW и изменить категорию у уже существующей карточки. На широком экране группы располагаются по две в строке, на узком --- по одной. В заголовке группы отображается число задач. Завершённые задачи остаются доступными в списке и на доске, но в представлении для планирования текущего объёма не показываются.

При переносе задачи из Must в Could сервер обновляет только \texttt{moscow}; статус, приоритет и срок сохраняются. Это важно для согласованности: решение о составе этапа не должно незаметно менять дату выполнения или результат матрицы. Для возврата в группу «Без категории» в PATCH передаётся \texttt{null}. Невозможность присвоить неизвестное значение, например \texttt{later}, проверяется моделью запроса и ограничением столбца SQLite.

Методика MoSCoW в прототипе реализована как персональный или общий для всех подключённых клиентов способ маркировки проектных задач. Система не знает, кто именно присвоил категорию, и не защищает её от изменения другим пользователем. Поэтому корректно говорить о подготовке структуры для последующей командной работы, но нельзя описывать текущую версию как систему назначения ролей, согласования решений или персональной ответственности исполнителей.

\subsection{Доска, список и календарь}

Помимо двух методов приоритизации, экран задач имеет представления «Доска», «Список» и «Календарь». Доска группирует задачи по состояниям: сделать, в работе, готово. Смена состояния отражает фактический ход выполнения и независима от важности, срочности и категории MoSCoW. В списке удобно искать задачи по заголовку, описанию и тегам, применять фильтры и сортировку. Такое сочетание позволяет сначала выбрать проект и вид, затем сузить набор данных до интересующих задач.

На рисунке~\ref{fig:tasks-board} показаны доска задач, панель поиска и фильтров, переключатель представлений и список проектов в текущей версии приложения для Windows.

\begin{figure}[!htbp]
\centering
\XeTeXpicfile figures/windows_board.png width \linewidth
\caption{Доска задач в приложении Windows}
\label{fig:tasks-board}
\end{figure}

Календарь получает с сервера задачи, у которых крайний срок приходится на выбранный месяц. Ответ группируется по дням. Записи расписания добавляются при включении отдельного параметра; их можно показывать и скрывать без изменения самих задач. На экране выбранный день открывает перечень соответствующих задач и мероприятий. Это уменьшает риск перепутать дедлайн задачи с датой повторяющейся встречи. Задача без дедлайна остаётся в общем списке, но не появляется в календаре крайнего срока.

Для просроченных задач есть отдельный режим отображения. Пользователь видит нарушение срока и может выбрать новое допустимое значение даты. При переносе дата начала остаётся нижней границей: новое окончание не должно оказаться раньше неё. Принцип неавтоматического изменения срока сохраняет ответственность пользователя за план; приложение обнаруживает проблему и предоставляет действие, но не назначает произвольную новую дату.

На уровне клиента загрузка данных сопровождается состояниями ожидания и ошибки. \texttt{ApiClient} ограничивает продолжительность HTTP-операции и переводит сбой соединения в понятное сообщение. После сохранения карточки экран повторно загружает данные. В \texttt{TasksPage} есть счётчики запросов, помогающие не подменять более свежий результат поздно завершившимся старым запросом. Такая деталь особенно актуальна при быстром переключении проектов и фильтров.

\subsection{Заметки и материалы задачи}

Заметка представляет собой самостоятельную сущность, а не длинное поле задачи. Пользователь может хранить поясняющий текст и связать его с несколькими задачами; при необходимости из заметки создаётся новая задача. Текст состоит из блоков, а блоки --- из фрагментов с атрибутами оформления. Модель допускает абзац, заголовок, маркированный и нумерованный список, три семейства шрифтов, размер текста, цвет и начертание. На сервере структура проверяется и сохраняется как JSON. Это позволяет менять оформление отдельного фрагмента без введения отдельной SQL-таблицы для каждого слова.

Связь заметки и задачи хранится в \texttt{note\_tasks}. При создании задачи из заметки заголовок и текст заметки используются как исходный материал, а затем между объектами создаётся связь. Заметка при этом не удаляется. Последующее изменение заметки не переписывает автоматически созданную задачу: у них разный жизненный цикл. В сценарии подготовки к посевной заметка с результатами осмотра может быть источником задачи «Заменить обнаруженный изношенный узел».

Изображения и аудио отправляются отдельным запросом к заметке. Сервер ограничивает допустимые типы и размер, а также сопоставляет заявленный тип с содержимым файла. В SQLite сохраняются название, MIME-тип, размер и внутреннее имя, а сами байты хранятся в каталоге загрузок. Такой способ не превращает текст заметки в контейнер большого бинарного файла. Для локального прототипа это достаточно просто, но в эксплуатации необходимо резервировать файлы одновременно с базой и ограничивать доступ пользователей к личным материалам.

Редактор заметок на Flutter, страницы списка и диалоги работы с файлами вынесены в отдельный каталог. Тем самым задачи и заметки имеют собственные компоненты интерфейса, но сохраняют возможность перехода друг к другу. Проверка медиа и форматированного содержимого выполняется на сервере, поскольку ограничение только в форме клиента не защищало бы API от прямого некорректного запроса.

\subsection{Повторяющееся расписание}

Расписание позволяет создать мероприятие на конкретную дату, каждую неделю или каждые две недели. Запись содержит название, описание, начальную дату, необязательную конечную дату, время начала и конца, а также связь с проектом. Сервер проверяет согласованность дат и интервала времени. В отличие от задачи, мероприятие обозначает событие в календаре; его повторение не создаёт новую задачу при каждом наступлении даты.

Метод получения экземпляров принимает начальную и конечную границы. Сервис находит записи, вычисляет даты повторения в пределах диапазона и исключает даты из \texttt{schedule\_exceptions}. Диапазон ограничен 366 днями, чтобы запрос не порождал произвольно большой результат. Отмена одной даты выполняется через маршрут исключения; восстановление отмены удаляет соответствующую строку исключения. Редактирование всей записи и удаление серии выполняются отдельными операциями.

Например, для контрольного осмотра оборудования раз в две недели выбирается начальная дата и значение \texttt{repeat\_every\_weeks\ =\ 2}. Если в одну из рассчитанных дат работы невозможны, пользователь отменяет только это событие. Следующие экземпляры продолжают рассчитываться от исходной даты, а не от даты исключения. Это поведение соответствует ожиданию от повторяющейся серии и экономит число хранимых строк.

В Flutter есть страница расписания и диалог его создания или редактирования. Из календаря можно открыть создание события на выбранную дату; отдельный переключатель добавляет занятия и мероприятия к календарным задачам. Если пользователь удаляет проект, записи расписания могут остаться без связи с ним, как и задачи. Такое поведение сохраняет данные, но требует от пользователя при необходимости повторно распределить их по проектам.

\subsection{Сквозной сценарий применения в АПК}

Для проверки предметной применимости рассмотрим учебную ситуацию: ответственному сотруднику нужно подготовить несколько работ до начала посевной кампании. Сначала он создаёт проект «Подготовка к посевной». В проект добавляются задачи с разными сроками и приоритетами. Затем к каждой задаче при необходимости относятся теги, заметки и мероприятия расписания. После этого пользователь открывает матрицу и MoSCoW. Такой сценарий проверяет не отдельную кнопку, а согласованность нескольких подсистем вокруг одного проекта.

Предположим, что текущая дата демонстрации --- 21 ноября 2026 года. Для задачи «Проверить исправность сеялки» указаны высокий приоритет, дедлайн 22 ноября и категория Must. По правилу матрицы это «срочно и важно»; в MoSCoW задача включена в обязательный объём. «Составить перечень запасных частей» с высоким приоритетом и сроком 30 ноября попадает в «важно, но не срочно» и может иметь категорию Should. «Уточнить время ближайшей встречи» с низким приоритетом и дедлайном 21 ноября относится к срочным, но менее важным, а в MoSCoW может оставаться без категории. «Оформить архив старых заметок» без близкого срока и с низким приоритетом оказывается в четвёртом секторе; её можно пометить Could или Won't в зависимости от состава текущего этапа. Даты и назначения в этом примере являются тестовыми, а не рекомендацией по организации реального хозяйства.

\begingroup\small
\begin{longtable}[]{@{}
  >{\raggedright\arraybackslash}p{(\linewidth - 8\tabcolsep) * \real{0.2000}}
  >{\raggedright\arraybackslash}p{(\linewidth - 8\tabcolsep) * \real{0.2000}}
  >{\raggedright\arraybackslash}p{(\linewidth - 8\tabcolsep) * \real{0.2000}}
  >{\raggedright\arraybackslash}p{(\linewidth - 8\tabcolsep) * \real{0.2000}}
  >{\raggedright\arraybackslash}p{(\linewidth - 8\tabcolsep) * \real{0.2000}}@{}}
\toprule\noalign{}
\begin{minipage}[b]{\linewidth}\raggedright
\textbf{Задача}
\end{minipage} & \begin{minipage}[b]{\linewidth}\raggedright
\textbf{Приоритет}
\end{minipage} & \begin{minipage}[b]{\linewidth}\raggedright
\textbf{Срок}
\end{minipage} & \begin{minipage}[b]{\linewidth}\raggedright
\textbf{Сектор матрицы}
\end{minipage} & \begin{minipage}[b]{\linewidth}\raggedright
\textbf{MoSCoW}
\end{minipage} \\
\midrule\noalign{}
\endhead
\bottomrule\noalign{}
\endlastfoot
Проверить исправность сеялки & high & 22.11.2026 & Срочно и важно & Must \\
Составить перечень запасных частей & high & 30.11.2026 & Важно, не срочно & Should \\
Уточнить время встречи & low & 21.11.2026 & Срочно, менее важно & Без категории \\
Оформить архив заметок & low & не задан & Не срочно, менее важно & Could \\
\end{longtable}
\endgroup

Таблица 7 --- Учебный набор задач для демонстрации двух методик

Пользователь связывает с первой задачей заметку «Осмотр оборудования», где фиксирует выявленные дефекты и добавляет фотографию детали. Если осмотр показывает необходимость отдельной работы, из заметки создаётся задача на ремонт. В расписании задаётся контрольное мероприятие на 22 ноября. Календарь показывает одновременно срок задачи и мероприятие в этот день, если показ расписания включён. Дальше задача переводится из состояния «Сделать» в «В работе» и после выполнения --- в «Готово». Она исчезает из матрицы и текущей MoSCoW-группировки, но сохраняется в истории списка задач.

Этот пример показывает, что пользовательское значение не исчерпывается автоматической классификацией. Проект формирует контекст, календарь помогает соблюсти время, заметка сохраняет сведения, а две группировки обеспечивают разные способы принятия решения. Если срок ремонта меняется, срочность задачи может измениться автоматически; категория Must останется прежней до ручной корректировки. Следовательно, один и тот же объект может участвовать в нескольких представлениях без противоречивого дублирования.

Сценарий имеет ограничение: приложение не рассчитывает последовательность технологических операций. Например, оно не запрещает отметить «Начать посев» выполненной до проверки оборудования, поскольку зависимости между задачами пока не моделируются. Для учебной версии это допускает гибкое использование универсального планировщика. Для более строгого управления процессом АПК следует добавить этапы, зависимость «предшественник --- последующая задача» и проверку готовности ресурсов.

\subsection{Методика и результаты проверки}

Проверка серверной части выполнена автоматическими тестами на временной SQLite-базе. Тестовый набор запускается командой \path{./.venv/Scripts/python.exe\ -m\ unittest\ discover\ -s\ tests\ -q} из корня проекта. На проверенном состоянии 21 ноября 2026 года выполнено \textbf{27 тестов, результат OK}. В тестах воспроизводятся операции создания, изменения, фильтрации и удаления задач, работа с заметками и расписанием, группировка матрицы и MoSCoW. Временный файл базы изолирует проверки от пользовательского файла \path{backend/tasks.db}.

Для матрицы тест создаёт задачи, покрывающие четыре сектора. Проверяется проектный фильтр, исключение выполненной задачи, вычисление срочности по дедлайну, явное отрицание срочности и возврат к автоматическому правилу при сбросе \texttt{urgent}. Дополнительно тест миграции создаёт старую схему и проверяет, что инициализация добавляет новый столбец без потери существующей задачи. Такой набор позволяет проверить и правило классификации, и совместимость с данными до обновления.

Для MoSCoW проверяется, что задачи выбранного проекта попадают в нужные группы, задачи другого проекта не смешиваются с ними, а завершённые не отображаются в текущем планировании. Перенос Must → Could не должен менять статус, приоритет и срок. Отдельно проверяются возврат в категорию без назначения, сохранение данных при повторном запуске и отклонение неизвестного кода. Эти случаи важнее простого подсчёта числа карточек, потому что подтверждают независимость методики от остальных свойств задачи.

\begingroup\small
\begin{longtable}[]{@{}
  >{\raggedright\arraybackslash}p{(\linewidth - 4\tabcolsep) * \real{0.3333}}
  >{\raggedright\arraybackslash}p{(\linewidth - 4\tabcolsep) * \real{0.3333}}
  >{\raggedright\arraybackslash}p{(\linewidth - 4\tabcolsep) * \real{0.3333}}@{}}
\toprule\noalign{}
\begin{minipage}[b]{\linewidth}\raggedright
\textbf{Проверяемое условие}
\end{minipage} & \begin{minipage}[b]{\linewidth}\raggedright
\textbf{Ожидаемый результат}
\end{minipage} & \begin{minipage}[b]{\linewidth}\raggedright
\textbf{Подтверждение}
\end{minipage} \\
\midrule\noalign{}
\endhead
\bottomrule\noalign{}
\endlastfoot
Срок завтра, высокий приоритет & Срочно и важно & Автоматический тест матрицы \\
Ручное \texttt{urgent\ =\ false} & Не срочно при близком сроке & Автоматический тест матрицы \\
Сброс \texttt{urgent\ =\ null} & Возврат к расчёту по сроку & Автоматический тест матрицы \\
Статус \texttt{done} & Отсутствие в активных группах & Тесты матрицы и MoSCoW \\
Смена \texttt{moscow} & Не меняются срок, статус и приоритет & Автоматический тест MoSCoW \\
Старые данные без новых колонок & Добавлены поля, задача сохранена & Тесты миграции \\
Неизвестная категория MoSCoW & Запрос отклоняется & Проверка модели \\
\end{longtable}
\endgroup

Таблица 8 --- Существенные проверки методов приоритизации

Отдельные тесты клиента находятся в \path{app/test}: там проверяются модели и вызовы API, заметки, расписание, матрица и MoSCoW. Статический анализ Dart ранее не выявил ошибок. При подготовке этой работы запуск \texttt{flutter\ test} в данной среде не был подтверждён: инструмент не смог надёжно завершить проверку интерфейса. Поэтому нельзя приравнивать наличие тестовых файлов к успешному выполнению всех Flutter-тестов. Полная демонстрация на целевых устройствах требует дополнительно запустить клиент и тесты в настроенном окружении пользователя.

Для проверки взаимодействия на уровне интерфейса после сборки достаточно выполнить последовательность: запустить сервер, открыть клиент, создать проект и четыре учебные задачи, проверить их положение в матрице, назначить и изменить категории MoSCoW, добавить заметку и событие расписания, затем повторно открыть проект. Результат каждой операции должен совпасть с серверными данными. Эта последовательность описывает программу ручной приёмки, а не заявление о завершённом испытании на Android в данной работе.

\subsection{Оценка решения и направления развития}

Разработанный прототип объединяет несколько способов просмотра одних и тех же задач без отдельного хранения карточек для каждого экрана. Это позволяет при изменении статуса или дедлайна получить актуальную запись и в списке, и в календаре, и в матрице. Правила матрицы формально определены, а ручное исключение помогает учесть обстоятельства, которые нельзя вывести из одной даты. MoSCoW предоставляет отдельный способ решать, какие работы необходимы на текущем этапе проекта. Проверки подтверждают корректность основных серверных переходов и сохранение ранее созданных данных при добавлении колонок.

Вместе с тем оценка не должна подменять пользовательское исследование. В этой работе не измерялись скорость выполнения задач, снижение числа просрочек, экономия средств хозяйства или точность экспертных решений. Демонстрационный набор показывает техническую применимость, но не доказывает улучшение производственных показателей. Для такого вывода потребовались бы пилотное использование, исходные и итоговые показатели, согласованные критерии оценки и наблюдение за достаточно длительным периодом работ.

Первая группа ограничений связана с пользователями. Нет регистрации, ролей, прав и истории изменений. Если несколько человек одновременно обращаются к одному серверу, они видят общий набор данных и могут менять категории друг друга. Для командной работы следует ввести пользователей, участие в проекте, назначенного исполнителя и проверку прав для каждого действия. Только после этого MoSCoW можно превратить в инструмент согласованного выбора состава работ с понятной ответственностью.

Вторая группа ограничений связана с производственным процессом. Задача пока не имеет зависимости от другой задачи, нормативной длительности, потребности в материале или технике. В универсальном планировщике эти детали можно описать в тексте заметки, но они не участвуют в автоматической проверке. Следующий этап --- модель этапов и зависимостей. Она позволит проверять предварительные условия и показывать, какие операции задерживаются из-за невыполненного предшественника. В агропромышленном сценарии отдельно может понадобиться учёт участков, техники и сезонных ограничений, однако соответствующие сущности следует вводить после анализа реальных пользователей, а не ради названия курсовой.

Третья группа ограничений относится к эксплуатации. Сервер хранит данные на локальном диске, а клиенту необходимо соединение с ним. Для устойчивой работы вне одного компьютера требуются безопасное сетевое размещение, HTTPS, резервное копирование SQLite и файловых вложений, контроль одновременного доступа, настройка часовых поясов и проверка Android-сборки на устройстве. Автономный режим с последующей синхронизацией был бы отдельной задачей: он требует локального хранилища на клиенте и правил разрешения конфликтов изменений.

Перспективным улучшением методов приоритизации является возможность задавать важность отдельно от обычного трёхуровневого приоритета и настраивать порог автоматической срочности для конкретного проекта. Сейчас срок в пределах двух дней --- единое правило для всех видов работ. Для отдельных сельскохозяйственных операций этот горизонт может быть слишком коротким либо слишком длинным. Изменяемый порог должен сопровождаться понятным объяснением пользователю, почему задача оказалась в данном секторе. Пока таких настроек в прототипе нет, поэтому в описании результата используется именно фиксированное правило.

\section*{Заключение}\addcontentsline{toc}{section}{Заключение}

В рамках курсовой работы разработан прототип кроссплатформенного приложения для управления проектами, задачами, заметками и расписанием. Для клиентской части использован Flutter, для HTTP API --- FastAPI, для хранения --- SQLite. Задачи можно создавать, изменять, фильтровать и показывать в нескольких представлениях. Календарь связывает сроки с днями месяца и при необходимости отображает события расписания. Заметки сохраняют контекст работы и могут быть связаны с задачами.

Главный механизм поддержки приоритизации --- матрица Эйзенхауэра. Она определяет важность через высокий приоритет, срочность --- через близкий дедлайн или явную отметку пользователя, после чего помещает активную задачу в один из четырёх секторов. MoSCoW дополняет матрицу ручным решением о необходимости задачи на текущем этапе проекта. Независимость двух классификаций позволяет одновременно оценивать и момент выполнения, и состав работ. Это решение основано на правилах с явной логикой; машинное обучение в приложении не реализовано.

Проверка актуального серверного состояния дала результат 27 успешных автоматизированных тестов. Они охватывают в том числе распределение по секторам, ручную корректировку срочности, смену категории MoSCoW и сохранение старых данных после изменения схемы. Успешное выполнение Flutter-тестов и полноценная работа на Android в рамках данной проверки не подтверждены. Пример «Подготовка к посевной» демонстрирует, как универсальный планировщик можно применить к набору работ АПК, не приписывая системе несуществующие отраслевые функции.

Развитие проекта предполагает добавление участников и прав, назначения исполнителей, зависимостей между задачами, устойчивого сетевого развёртывания и испытаний на целевых устройствах. Эти изменения позволят приблизиться от персонального планирования к управлению реальными производственными процессами. На текущем этапе получен работоспособный программный прототип с воспроизводимыми правилами приоритизации и ясными границами применения.

\section*{Список использованных источников}\addcontentsline{toc}{section}{Список использованных источников}

\begin{enumerate}
\def\labelenumi{\arabic{enumi}.}
\tightlist
\item
  Flutter. Desktop support for Flutter {[}Электронный ресурс{]}. URL: \url{https://docs.flutter.dev/platform-integration/desktop} (дата обращения: 21.11.2026).
\item
  FastAPI. Bigger Applications --- Multiple Files {[}Электронный ресурс{]}. URL: \url{https://fastapi.tiangolo.com/tutorial/bigger-applications/} (дата обращения: 21.11.2026).
\item
  FastAPI. Features {[}Электронный ресурс{]}. URL: \url{https://fastapi.tiangolo.com/features/} (дата обращения: 21.11.2026).
\item
  SQLite. Foreign Key Support {[}Электронный ресурс{]}. URL: \url{https://www.sqlite.org/foreignkeys.html} (дата обращения: 21.11.2026).
\item
  Agile Business Consortium. What is MoSCoW Prioritization? {[}Электронный ресурс{]}. URL: \url{https://www.agilebusiness.org/resource/what-is-moscow-prioritization/} (дата обращения: 21.11.2026).
\item
  Flutter. Testing Flutter apps {[}Электронный ресурс{]}. URL: \url{https://docs.flutter.dev/testing/overview} (дата обращения: 21.11.2026).
\item
  Python Software Foundation. sqlite3 --- DB-API 2.0 interface for SQLite databases {[}Электронный ресурс{]}. URL: \url{https://docs.python.org/3/library/sqlite3.html} (дата обращения: 21.11.2026).
\item
  Исходный код проекта WorkNotes: модули \texttt{backend}, \texttt{app}, тесты \texttt{tests} и \path{app/test}, документация \texttt{docs} (локальный репозиторий, состояние на 21.11.2026).
\end{enumerate}

\end{document}
